% =========================================================
% Quantifier Laws
% =========================================================

\subsection*{Quantifier Laws}

\begin{remark*}[Section toolkit: Quantifier algebra]
Quantifier algebra is the predicate-logic analogue of propositional
equivalence laws. These laws justify rewriting quantified formulas while
preserving logical equivalence.
\end{remark*}

\begin{theorembox}{Theorem (Quantifier Negation Laws)}
\begin{theorem}[Quantifier Negation Laws]
\label{thm:qneg}
For every formula \(\varphi\),
\[
\neg\forall x\,\varphi \equiv \exists x\,\neg\varphi,
\qquad
\neg\exists x\,\varphi \equiv \forall x\,\neg\varphi.
\]
\hyperref[prf:qneg]{Go to proof.}
\end{theorem}
\end{theorembox}

\begin{remark*}[Standard quantified statement]
\[
\forall \varphi\in\mathsf{Form}_{\mathcal L}\quad
\neg\forall x\,\varphi \equiv \exists x\,\neg\varphi
\quad\text{and}\quad
\neg\exists x\,\varphi \equiv \forall x\,\neg\varphi.
\]
\end{remark*}

<<<<<<< HEAD
\begin{example*}[Negating a universal implication]
For a universal class statement, the negation says that there is a
counterexample:
\[
\neg\bigl(\forall x\,(P(x)\rightarrow Q(x))\bigr)
\equiv
\exists x\,\neg(P(x)\rightarrow Q(x))
\equiv
\exists x\,(P(x)\wedge \neg Q(x)).
\]
Thus ``not every $P$ is a $Q$'' means ``some object is a $P$ and not a $Q$.''
\end{example*}

\begin{example*}[Negating an existential conjunction]
For an existential conjunction, the negation says that no object satisfies both
conditions:
\[
\neg\bigl(\exists x\,(P(x)\wedge Q(x))\bigr)
\equiv
\forall x\,\neg(P(x)\wedge Q(x))
\equiv
\forall x\,(\neg P(x)\vee \neg Q(x))
\equiv
\forall x\,(P(x)\rightarrow \neg Q(x)).
\]
Thus ``there is no object that is both $P$ and $Q$'' may be read as ``every
$P$ fails to be a $Q$.''
\end{example*}

\begin{remark}[Common error]
Negating only the predicate without flipping the quantifier is incorrect:
$\neg(\forall x\,P(x)) \not\equiv \forall x\,\neg P(x)$.
The right side asserts that \emph{every} element fails $P$, which is much
stronger than merely asserting that not every element satisfies $P$.
\end{remark}
=======
\begin{remark*}[Predicate reading]
\(\operatorname{LogicalEquivalence}(\neg\forall x\,\varphi,\exists x\,\neg\varphi)\)
and
\(\operatorname{LogicalEquivalence}(\neg\exists x\,\varphi,\forall x\,\neg\varphi)\).
\end{remark*}
>>>>>>> 591349ed2a1982198e862953b73959c91a638d59

\begin{remark*}[Interpretation]
Negating a quantified statement flips the quantifier and pushes the negation
into the scope.
\end{remark*}

\begin{dependencies}
\begin{itemize}
  \item \hyperref[def:universal-q]{Universal Formula}
  \item \hyperref[def:existential-q]{Existential Formula}
  \item \hyperref[def:log-equiv]{Logical Equivalence}
\end{itemize}
\end{dependencies}

\begin{theorembox}{Theorem (Double Negation and Quantifiers)}
\begin{theorem}[Double Negation and Quantifiers]
\label{thm:quantifier-double-negation}
For every formula \(\varphi\),
\[
\neg\neg\forall x\,\varphi \equiv \forall x\,\varphi,
\qquad
\neg\neg\exists x\,\varphi \equiv \exists x\,\varphi.
\]
\hyperref[prf:quantifier-double-negation]{Go to proof.}
\end{theorem}
\end{theorembox}

\begin{remark*}[Standard quantified statement]
\[
\forall \varphi\in\mathsf{Form}_{\mathcal L}\quad
\neg\neg\forall x\,\varphi \equiv \forall x\,\varphi
\quad\text{and}\quad
\neg\neg\exists x\,\varphi \equiv \exists x\,\varphi.
\]
\end{remark*}

\begin{remark*}[Predicate reading]
\(\operatorname{LogicalEquivalence}(\neg\neg Qx\,\varphi,Qx\,\varphi)\) for
\(Q\in\{\forall,\exists\}\).
\end{remark*}

\begin{remark*}[Interpretation]
Double negation does not change the quantifier form once the two negations are
cancelled.
\end{remark*}

\begin{dependencies}
\begin{itemize}
  \item \hyperref[thm:qneg]{Quantifier Negation Laws}
\end{itemize}
\end{dependencies}

\begin{theorembox}{Theorem (Same-Type Quantifier Commutation)}
\begin{theorem}[Same-Type Quantifier Commutation]
\label{thm:qcomm}
Quantifiers of the same type commute:
\[
\forall x\,\forall y\,\varphi \equiv \forall y\,\forall x\,\varphi,
\qquad
\exists x\,\exists y\,\varphi \equiv \exists y\,\exists x\,\varphi.
\]
\hyperref[prf:qcomm]{Go to proof.}
\end{theorem}
\end{theorembox}

\begin{remark*}[Standard quantified statement]
\[
\forall \varphi\in\mathsf{Form}_{\mathcal L}\quad
\forall x\,\forall y\,\varphi \equiv \forall y\,\forall x\,\varphi
\quad\text{and}\quad
\exists x\,\exists y\,\varphi \equiv \exists y\,\exists x\,\varphi.
\]
\end{remark*}

\begin{remark*}[Predicate reading]
The theorem asserts logical equivalence after swapping adjacent quantifiers of
the same type.
\end{remark*}

\begin{remark*}[Interpretation]
Two universal choices can be checked in either order, and two existential
witnesses can be selected in either order.
\end{remark*}

\begin{dependencies}
\begin{itemize}
  \item \hyperref[def:universal-q]{Universal Formula}
  \item \hyperref[def:existential-q]{Existential Formula}
\end{itemize}
\end{dependencies}

\begin{theorembox}{Theorem (Mixed Quantifiers Do Not Commute)}
\begin{theorem}[Mixed Quantifiers Do Not Commute]
\label{thm:mixed-quantifiers-do-not-commute}
In general,
\[
\forall x\,\exists y\,\varphi
\not\equiv
\exists y\,\forall x\,\varphi.
\]
\hyperref[prf:mixed-quantifiers-do-not-commute]{Go to proof.}
\end{theorem}
\end{theorembox}

\begin{remark*}[Standard quantified statement]
\[
\exists \varphi\in\mathsf{Form}_{\mathcal L}\quad
\forall x\,\exists y\,\varphi
\not\equiv
\exists y\,\forall x\,\varphi.
\]
\end{remark*}

\begin{remark*}[Predicate reading]
\(\neg\operatorname{LogicalEquivalence}(\forall x\,\exists y\,\varphi,
\exists y\,\forall x\,\varphi)\) in general.
\end{remark*}

\begin{remark*}[Interpretation]
The first form permits the witness \(y\) to depend on \(x\). The second form
requires one uniform witness \(y\) that works for every \(x\).
\end{remark*}

\begin{dependencies}
\begin{itemize}
  \item \hyperref[thm:qcomm]{Same-Type Quantifier Commutation}
\end{itemize}
\end{dependencies}

\begin{theorembox}{Theorem (Valid Quantifier Distribution)}
\begin{theorem}[Valid Quantifier Distribution]
\label{thm:qdist}
For all formulas \(\varphi,\psi\),
\[
\forall x\,(\varphi\land\psi)
\equiv
(\forall x\,\varphi)\land(\forall x\,\psi),
\]
\[
\exists x\,(\varphi\lor\psi)
\equiv
(\exists x\,\varphi)\lor(\exists x\,\psi).
\]
\hyperref[prf:qdist]{Go to proof.}
\end{theorem}
\end{theorembox}

\begin{remark*}[Standard quantified statement]
\[
\forall \varphi,\psi\in\mathsf{Form}_{\mathcal L}\quad
\forall x\,(\varphi\land\psi)
\equiv
(\forall x\,\varphi)\land(\forall x\,\psi)
\]
and
\[
\exists x\,(\varphi\lor\psi)
\equiv
(\exists x\,\varphi)\lor(\exists x\,\psi).
\]
\end{remark*}

\begin{remark*}[Predicate reading]
The displayed universal-conjunction and existential-disjunction distributions
are logical equivalences.
\end{remark*}

\begin{remark*}[Interpretation]
Universal quantification distributes over conjunction because both conjuncts
must hold for every value. Existential quantification distributes over
disjunction because a witness for either disjunct witnesses the disjunction.
\end{remark*}

\begin{dependencies}
\begin{itemize}
  \item \hyperref[def:log-equiv]{Logical Equivalence}
  \item \hyperref[def:universal-q]{Universal Formula}
  \item \hyperref[def:existential-q]{Existential Formula}
\end{itemize}
\end{dependencies}

\begin{theorembox}{Theorem (Invalid Quantifier Distribution)}
\begin{theorem}[Invalid Quantifier Distribution]
\label{thm:invalid-quantifier-distribution}
In general,
\[
\forall x\,(\varphi\lor\psi)
\not\equiv
(\forall x\,\varphi)\lor(\forall x\,\psi),
\]
\[
\exists x\,(\varphi\land\psi)
\not\equiv
(\exists x\,\varphi)\land(\exists x\,\psi).
\]
\hyperref[prf:invalid-quantifier-distribution]{Go to proof.}
\end{theorem}
\end{theorembox}

\begin{remark*}[Standard quantified statement]
\[
\exists \varphi,\psi\in\mathsf{Form}_{\mathcal L}\quad
\forall x\,(\varphi\lor\psi)
\not\equiv
(\forall x\,\varphi)\lor(\forall x\,\psi),
\]
\[
\exists \varphi,\psi\in\mathsf{Form}_{\mathcal L}\quad
\exists x\,(\varphi\land\psi)
\not\equiv
(\exists x\,\varphi)\land(\exists x\,\psi).
\]
\end{remark*}

\begin{remark*}[Predicate reading]
The displayed distribution patterns are not logical equivalences in general.
\end{remark*}

\begin{remark*}[Interpretation]
The invalid laws fail because witnesses or counterexamples may vary between
the two subformulas.
\end{remark*}

\begin{dependencies}
\begin{itemize}
  \item \hyperref[thm:qdist]{Valid Quantifier Distribution}
\end{itemize}
\end{dependencies}

\begin{theorembox}{Theorem (Vacuous Quantification)}
\begin{theorem}[Vacuous Quantification]
\label{thm:vacuous}
If \(x\) does not occur free in \(\varphi\), then
\[
\forall x\,\varphi\equiv\varphi,
\qquad
\exists x\,\varphi\equiv\varphi.
\]
\hyperref[prf:vacuous]{Go to proof.}
\end{theorem}
\end{theorembox}

\begin{remark*}[Standard quantified statement]
\[
x\notin\mathsf{FV}(\varphi)
\Rightarrow
(\forall x\,\varphi\equiv\varphi)
\land
(\exists x\,\varphi\equiv\varphi).
\]
\end{remark*}

\begin{remark*}[Predicate reading]
Vacuous quantification preserves logical equivalence when the quantified
variable is not free in the scope.
\end{remark*}

\begin{remark*}[Interpretation]
A quantifier that binds no occurrence has no semantic effect.
\end{remark*}

\begin{dependencies}
\begin{itemize}
  \item \hyperref[def:quantifier-scope]{Quantifier Scope}
\end{itemize}
\end{dependencies}

\begin{theorembox}{Theorem (Renaming Bound Variables)}
\begin{theorem}[Renaming Bound Variables]
\label{thm:rename}
If \(y\) is free for \(x\) in \(\varphi\) and no capture is introduced, then
\[
\forall x\,\varphi\equiv\forall y\,\varphi[y/x],
\qquad
\exists x\,\varphi\equiv\exists y\,\varphi[y/x].
\]
\hyperref[prf:rename]{Go to proof.}
\end{theorem}
\end{theorembox}

\begin{remark*}[Standard quantified statement]
\[
y \text{ free for } x \text{ in } \varphi
\Rightarrow
(\forall x\,\varphi\equiv\forall y\,\varphi[y/x])
\land
(\exists x\,\varphi\equiv\exists y\,\varphi[y/x]).
\]
\end{remark*}

\begin{remark*}[Predicate reading]
Renaming bound variables preserves logical equivalence when the renaming is
capture-avoiding.
\end{remark*}

\begin{remark*}[Interpretation]
Bound variable names are placeholders. Renaming them changes notation, not
meaning, provided no free variable becomes captured.
\end{remark*}

\begin{dependencies}
\begin{itemize}
  \item \hyperref[def:free-for]{Free for Substitution}
  \item \hyperref[def:subst-notation]{Substitution Notation}
\end{itemize}
\end{dependencies}
